\documentclass[10pt,a4paper]{amsart}
\usepackage[margin=19mm]{geometry}
\usepackage{amsmath,amssymb,mathtools}
\usepackage[scr=rsfs,cal=euler]{mathalfa}
\usepackage{xfrac}
\usepackage[unicode,hidelinks,bookmarks=false]{hyperref}
\newcommand{\ie}{i.e.\ }
\newcommand{\cf}{cf.\ }
\newcommand{\resp}{respectively\ }
\newcommand{\Subsection}{\subsection}
\providecommand{\nobreakdash}{\nobreak\hbox{-}}
\setlength{\emergencystretch}{2em}
\multlinegap0pt
\usepackage{amssymb}
\usepackage{float}
\newtheorem{lemma}{Lemma}
\usepackage{cleveref}
\crefname{lemma}{Lemma}{Lemmas}
\providecommand{\Bigabs}[1]{\Bigl\lvert#1\Bigr\rvert}
\newcommand{\C}{\mathbb{C}}
\newcommand{\Htran}{\mathsf{H}}
\providecommand{\abs}[1]{\lvert#1\rvert}
\providecommand{\bignorm}[1]{\bigl\lVert#1\bigr\rVert}
\newcommand{\defi}{:=}
\newcommand{\kval}{\kappa_{\mathsf{val}}}
\newcommand{\kvec}{\kappa_{\mathsf{vec}}}
\providecommand{\norm}[1]{\lVert#1\rVert}
\newcommand{\range}{\mathsf{range}}
\providecommand{\spa}[1]{\mathsf{span}\{#1\}}
\title{Stabilizing the Rayleigh--Ritz procedure by randomization}
\author{Nian Shao}
\date{Source: arXiv:2604.01037v3 (CC BY 4.0)}
\begin{document}
\maketitle

\noindent\emph{Source and licence.} Stabilizing the Rayleigh--Ritz procedure by randomization. Version arXiv:2604.01037v3 (CC BY 4.0), distributed by arXiv under the Creative Commons Attribution 4.0 International licence, http://creativecommons.org/licenses/by/4.0/, as recorded in the arXiv licence metadata for this exact version and shown on the abstract page https://arxiv.org/abs/2604.01037v3. Full licence notice: \texttt{licenses/arxiv-2604.01037-CC-BY-4.0.txt}.\par\smallskip

\noindent\textit{Editorial scope.} Counted unit: the article's lemma lem:kappa (source0.tex lines 508--527). The article's complete original proof of it is delivered with it (source0.tex lines 529--577, one \texttt{proof} environment of 49 lines). No mathematics is written, repaired or re-derived: the statement and the proof are the article's own lines, in the article's own notation, and no retained line is substituted or re-flowed.  Retained uncounted as the source-local context the proof needs: lines 457--461; lines 492--503.  Nothing was omitted from any retained span, so the package carries no omission record.  No retained line contains a citation, so the wrapper carries no bibliography and no bibliography entry.  No retained line contains a figure environment, an image-inclusion command or drawing code, so no image is delivered and the package declares no asset dependency.  Editorial formatting only, in addition: an AMS \texttt{amsart} preamble that restates the article's own declarations and macros used by the retained lines, namely the environments \texttt{lemma} on separate counters, together with the standard packages those lines need.  The retained environments print in this document as Lemma 1 in order of appearance, whereas the article prints the same items with the numbering of its own full document.  The complete native source of the article as distributed in the arXiv e-print for this exact version is bundled unchanged as \texttt{sources/197631/source0.tex}.
\begin{lemma}
    \label{lem:exact}
    Let $(\lambda,v)$ be a simple eigenpair of $A(\cdot)$.  Given an arbitrary fixed full-rank matrix $W\in\C^{n\times m}$ satisfying $We_{1}=v$, the following statements hold for almost every $\Omega\in\C^{n\times m}$ with respect to Lebesgue measure:
    The matrix-valued function $B(\xi)=\Omega^{\Htran}A(\xi)W$ is regular, and $(\lambda,e_{1})$ is a simple eigenpair of $B(\cdot)$.
\end{lemma}

\begin{lemma}
    \label{lem:Gaussian}
    Let $\Omega\in\C^{n\times m}$ be a complex Gaussian random matrix. Then 
    \begin{equation*}
        \Pr\bigl(\norm{\Omega}\geq \sqrt{n}+\sqrt{m}+\sqrt{\ln(2/\delta)}\bigr)\leq \delta
        \quad\text{for all}\quad \delta\geq 0. 
    \end{equation*}
    When $m=n$, that is, $\Omega\in\C^{m\times m}$ is a square matrix (or a scalar), it holds that:
    \begin{equation*}
        \Pr\bigl(\norm{\Omega^{-1}}\geq \sqrt{m/\delta}\bigr)\leq \delta \quad\text{for all}\quad \delta \geq 0. 
    \end{equation*}
\end{lemma}

\begin{lemma}
    \label{lem:kappa}
    Let $u$ and $v$ be unit left and right eigenvectors of $A(\cdot)$ corresponding to a simple eigenvalue $\lambda$.
    Given an arbitrary fixed orthonormal matrix $W$ admitting $W=[v,W_{\perp}]\in\C^{n\times m}$, we define $B(\xi)=\Omega^{\Htran}A(\xi)W$, where $\Omega\in\C^{n\times m}$ is a complex Gaussian random matrix. Then 
    \begin{equation*}
        \begin{aligned}
                \kval(B,\lambda) &\leq  \frac{1}{\sqrt{\delta}}\cdot \kval(A,\lambda) &\text{holds with probability at least $1-\delta$,}\\ 
                \kvec(B,\lambda) &\leq  \frac{\sqrt{m-1}}{\sqrt{\delta}}\cdot \kvec(A,\lambda) &\text{holds with probability at least $1-\delta$.}\\ 
        \end{aligned}
    \end{equation*}
    Here, $\kval(A,\lambda)$ and $\kvec(A,\lambda)$ are eigenvalue and eigenvector condition numbers of $A(\cdot)$:
    \begin{equation}
        \label{eq:defkappaA}
        \kval(A,\lambda)\defi \frac{\norm{u}\norm{v}}{\abs{u^{\Htran}A^{\prime}(\lambda)v}}
        \quad\text{and}\quad 
        \kvec(A,\lambda)\defi \bignorm{(Q_{\perp}^{\Htran}A(\lambda)V_{\perp})^{-1}},
    \end{equation}
    where $Q_{\perp}$ and $V_{\perp}$ are orthonormal bases of $\spa{A^{\prime}(\lambda)v}^{\perp}$ and $\spa{v}^{\perp}$, respectively.

\end{lemma}

\begin{proof}
    By \cref{lem:exact}, we know that $(\lambda,e_{1})$ is a simple eigenpair of $B(\cdot)$ almost surely. We further denote its corresponding unit left eigenvector by $x$.
    
    \emph{Eigenvalue condition number:}
    Since $(0,v)$ is a simple eigenpair of the matrix $A(\lambda)$ and $[v,W_{\perp}]$ is orthonormal, we can let $P\in\C^{n\times (m-1)}$ be an orthonormal basis of $\range(A(\lambda)W_{\perp})$ and $p_{\perp} = (I-PP^{\Htran})A^{\prime}(\lambda)v$. Since $u^{\Htran}A(\lambda)=0$, we know that $u^{\Htran}P=0$ and 
    \begin{equation*}
        \norm{p_{\perp}} \geq \norm{uu^{\Htran}(I-PP^{\Htran})A^{\prime}(\lambda)v}= \norm{uu^{\Htran}A^{\prime}(\lambda)v} =  \abs{u^{\Htran}A^{\prime}(\lambda)v} = \frac{1}{\kval(A,\lambda)}.
    \end{equation*}
    Since the eigenvalue $0$ of the matrix $B(\lambda)$ is also simple, we know that  
    \begin{equation*}
        \spa{x} = \range\bigl(B(\lambda)\bigr)^{\perp}=\range\bigl(\Omega^{\Htran}A(\lambda)W_{\perp}\bigr)^{\perp}=\range(\Omega^{\Htran}P)^{\perp},    
    \end{equation*}
    and in turn, $x^{\Htran}\Omega^{\Htran} p_{\perp} = x^{\Htran}\Omega^{\Htran} A^{\prime}(\lambda)v-x^{\Htran}\Omega^{\Htran} PP^{\Htran}A^{\prime}(\lambda)v = x^{\Htran}\Omega^{\Htran} A^{\prime}(\lambda)v$.
    Note that $\Omega^{\Htran}P$ and $\Omega^{\Htran}p_{\perp}$ are independent; consequently, $x$ and $\Omega^{\Htran}p_{\perp}$ are also independent. Thus, 
    \begin{equation*}
        \frac{1}{\kval(B,\lambda)\cdot\norm{p_{\perp}}} = \frac{\abs{x^{\Htran}B^{\prime}(\lambda)e_{1}}}{\norm{p_{\perp}}} = \frac{\abs{x^{\Htran}\Omega^{\Htran}A^{\prime}(\lambda)v}}{\norm{p_{\perp}}} = \frac{\abs{x^{\Htran}\Omega^{\Htran}p_{\perp}}}{\norm{p_{\perp}}} = \Bigabs{x^{\Htran}\Omega^{\Htran}\bigl(p_{\perp}/\norm{p_{\perp}}\bigr)}
    \end{equation*}
    is the magnitude of a complex Gaussian random variable. Using the tail bound in \cref{lem:Gaussian}, we know that, with probability at least $1-\delta$, it holds that 
    \begin{equation*}
        \kval(B,\lambda)\leq  \frac{1}{\sqrt{\delta}\cdot\norm{p_{\perp}}}  \leq \frac{1}{\sqrt{\delta}}\cdot\kval(A,\lambda).
    \end{equation*}

    \emph{Eigenvector condition number:} Let 
    \begin{equation*}
        Q_{\perp}Q_{\perp}^{\Htran}A(\lambda)W_{\perp} = QR
    \end{equation*}
    be the thin QR decomposition. Since the matrix $[A^{\prime}(\lambda)v,A(\lambda)V_{\perp}]$ is nonsingular and $\range(W_{\perp})\subset \range(V_{\perp})$. We know that $Q\in\C^{n\times (m-1)}$ and
    \begin{equation*}
        \norm{R^{-1}} \leq \bignorm{\bigl(Q_{\perp}^{\Htran}A(\lambda)V_{\perp}\bigr)^{-1}} = \kvec(A,\lambda),
    \end{equation*}
    where we use the interlacing property on the smallest singular values of $Q_{\perp}^{\Htran}A(\lambda)W_{\perp}$ and $Q_{\perp}^{\Htran}A(\lambda)V_{\perp}$. 
    Recall that 
    \begin{equation*}
        \kvec(B,\lambda) =  \bignorm{\bigl((\Omega Z_{\perp})^{\Htran}A(\lambda)W_{\perp}\bigr)^{-1}}, 
    \end{equation*}
    where $Z_{\perp}$ is an orthonormal basis of $\spa{\Omega^{\Htran}A^{\prime}(\lambda)v}^{\perp}$. We know that $(\Omega Z_{\perp})^{\Htran}A^{\prime}(\lambda)v=0$, hence $\Omega Z_{\perp} = Q_{\perp}Q_{\perp}^{\Htran}\Omega Z_{\perp}$. Thus, 
    \begin{equation*}
        \begin{aligned}
            \kvec(B,\lambda) &= 
        \bignorm{\bigl((\Omega Z_{\perp})^{\Htran}Q_{\perp}Q_{\perp}^{\Htran}A(\lambda)W_{\perp}\bigr)^{-1}} = \bignorm{\bigl((\Omega Z_{\perp})^{\Htran}QR\bigr)^{-1}} \\ 
        &\leq \bignorm{\bigl(Z_{\perp}^{\Htran}(\Omega^{\Htran}Q)\bigr)^{-1}}\cdot \norm{R^{-1}}\leq \bignorm{\bigl(Z_{\perp}^{\Htran}(\Omega^{\Htran}Q)\bigr)^{-1}}\cdot \kvec(A,\lambda).        
        \end{aligned}
    \end{equation*}
    Note that $A^{\prime}(\lambda)v$ and $Q$ are orthogonal. We know that $\Omega^{\Htran}A^{\prime}(\lambda)v$ and $\Omega^{\Htran}Q$ are independent; consequently, $Z_{\perp}$ and $\Omega^{\Htran}Q$ are also independent. Then $Z_{\perp}^{\Htran}(\Omega^{\Htran}Q)\in\C^{(m-1)\times (m-1)}$ is a complex Gaussian random matrix. By \cref{lem:Gaussian},  
    \begin{equation*}
        \bignorm{\bigl(Z_{\perp}^{\Htran}(\Omega^{\Htran}Q)\bigr)^{-1}}\leq \sqrt{(m-1)/\delta}
    \end{equation*} 
    holds with probability at least $1-\delta$. 
\end{proof}

\end{document}
